\section{Summary}

To sum up, the presented methods for measuring the impact of pruning on OOD and adversarial robustness of a neural network
showed that the proposed gradual magnitude-based pruning strategy does not reshape the model's robustness features. Evaluation on the CIFAR-10-C dataset
indicated that the pruned models did not get more vulnerable to common corruptions than the original dense model. Furthermore, when examining the progression
of the Fourier sensitivity heatmaps during the pruning process, no significant change could be detected in the model's robustness to 
perturbations of different frequencies and orientations, except at extreme sparsity levels. However, the distance to the decision boundary
- approximated by the norm of the minimal adversarial perturbation found by the FAB attack - gradually decreased as sparsity levels increased. Based on these
results, we hypothesize that our pruning process only decreases robustness to adversarial inputs, while maintains the OOD robustness of the model.

While the proposed pruning method did not worsen the OOD robustness, research suggests that model compression techniques like pruning
can be enhanced to contribute to robustness targets as well during training. Therefore, future work is needed to find modifications of this specific
pruning method to improve upon robustness, accuracy and compactness simultaneously.   